\section{Related works}
\subsection{Video Generation}

Recent advances in generative learning methods—including autoregressive models \cite{var}, diffusion models \cite{ddnm, ddim, im-ddpm}, and flow matching \cite{flowmatching}—have revolutionized video generation. Early efforts like the Video Diffusion Model \cite{video-dm} pioneered pixel-space denoising, while later works (Make-A-Video \cite{make-a-video}, PYoCo \cite{pyoco}, Imagen Video \cite{imagen-video}) integrated large language models for text-to-video synthesis. To tackle video dimensionality, research shifted toward latent space learning: VideoGPT \cite{VideoGPT} combined VQ-VAE \cite{vq-gan} with transformers, establishing foundational latent modeling. Subsequent innovations employed diffusion models to approximate latent distributions, spawning latent video diffusion frameworks (\cite{lvdm, magicvideo, make-your-vid, video-ldm, stable-video-diffusion, wang2023lavie, chen2023videocrafter1, chen2024videocrafter2}). Among these, CogVideoX \cite{yang2024cogvideox} introduced diffusion transformers and temporal-compressed VAEs, enhancing motion complexity. Leveraging expanded datasets and compute, modern large-scale generators (\cite{genmo2024mochi, kong2024hunyuanvideo, wan}) now achieve unprecedented quality.

\subsection{Audio-driven Human Animation}

Audio-driven human animation aims to generate dynamic videos from a static reference image, producing synchronized facial expressions and body movements based on audio control signals. In recent years, with the success of diffusion models, end-to-end audio-to-video synthesis methods \cite{tian2024emo,wei2024aniportrait,xu2024hallo,chen2025echomimic,cui2024hallo3,ji2024sonic,li2024latentsync,jiang2024loopy} have demonstrated considerable potential. These approaches eliminate the need for intermediate representations and exhibit superior performance in portrait animation for talking head generation. Another line of research extends talking head generation to talking body generation. Methods in this category \cite{lincyberhost,tian2025emo2,lin2025omnihuman,meng2024echomimicv2,gan2025omniavatar,wang2025fantasytalking} have achieved significant progress, demonstrating enhanced naturalness and consistent portrait animation capabilities, largely attributed to large-scale, high-quality training data. Recently, some works \cite{chen2025hunyuanvideo,kong2025let,huang2025bind} have attempted to advance from single human animation to multi-human animation.

The aforementioned human animation approaches have achieved satisfactory results in synthesizing short videos. However, when generating longer video sequences, they encounter error accumulation issues \cite{kong2025dam}, such as identity degradation and color deviations. A concurrent work, named StableAvatar \cite{stableavatar}, also achieves infinite-long sequence human animation. But it uses only a single image as the condition, and is not capable for the video-to-video task.
